% §6.3 LEACE + IGBP deconfounding of BERT embeddings, 12-city panel.
% Numbers locked against results/leace_12city/{city}_leace.json for all 12.

\subsection{Spatial deconfounding of BERT embeddings via LEACE + IGBP}
\label{sec:leace-12}

The sentence-BERT representation that underlies the Baur channel in
Section~\ref{sec:baur-12city} encodes more than the language a listing
agent uses. Pre-trained on a generic corpus, the encoder has no
particular incentive to separate semantic content about a property's
condition or amenities from geographic information that correlates
with the listing's neighborhood. If the embedding leaks
location-targeted information into the leading principal component,
the per-$\sigma$ effect we report in Section~\ref{sec:baur-12city}
inherits a spatial confounder mediated through the encoder, and our
identification claim is correspondingly weaker.

We address this confounder in two stages. The first stage applies
the LEACE projection of \citet{belrose2023leace} to the BERT
embedding, which projects out the linear subspace identified by the
mean-conditional-covariance estimator of the concept variable, where
the concept here is the lat-lon coordinate pair of the listing.
LEACE attains exact linear guardedness in the sense that no linear
classifier of the residual embedding can predict the concept better
than chance. We verify this attainment empirically: across the
twelve metropolitan markets in our panel, the post-LEACE
guardedness residual is $5\text{e-}17$ on average, ranging from
$3.7\text{e-}17$ to $8.9\text{e-}17$, at the numerical floor of
double-precision arithmetic. Linear concept erasure is complete in
every market.

The second stage applies the Iterative Gradient-Based Probing
correction of \citet{iskander2023log_linear}, which addresses a known
limitation of LEACE. Because LEACE projects out only the linear
subspace, the residual embedding can still encode the concept
non-linearly through an MLP probe. IGBP iterates between fitting an
MLP probe of the concept on the residual and projecting out the
direction in which the probe accumulates predictive mass. With five
outer iterations, IGBP reduces the adversarial probe's mean squared
error by an average of $17\%$ across the twelve cities, with the
reduction varying considerably by market: Chicago and Denver close
$71\%$ and $21\%$ of the residual nonlinear leakage respectively,
while Portland and Seattle exhibit essentially no residual nonlinear
leakage to close ($0.2\%$ reduction in both). The variation
across cities tracks the structural property of the data: markets
in which the listing language has stronger ZIP-targeted patterns
have more residual nonlinear leakage for IGBP to close, while
markets with more diffuse language patterns are already
near-deconfounded after LEACE alone.

\paragraph{Routing.} The LEACE machinery requires the analyst to
specify the concept as either a categorical variable (e.g., ZIP code
identifier) or a continuous one (e.g., the lat-lon coordinate pair).
Categorical LEACE has known instability properties when the
categorical cardinality is large relative to the sample size. Our
panel ranges from $24$ ZIP codes in San Francisco to $126$ in New
York, against per-city test sets of $87$--$105$ listings. We
therefore route the concept through the continuous lat-lon
representation in all twelve markets, which sidesteps the
categorical-cardinality issue without losing the spatial information
the projection is designed to remove.

\paragraph{Power.} For each market we report the binomial-power
diagnostic for the post-LEACE+IGBP linear-guardedness test on the
held-out fold. The test asks whether a linear classifier of the
residual embedding can predict the ZIP-targeted concept better than
chance, where chance is $1/K$ at $K$ ZIP codes per market. Across
the twelve markets the minimum detectable gain at $80\%$ power
against $\alpha = 0.05$ is $7\%$ on average above chance, ranging
from $4\%$ in New York (where the chance rate is $0.8\%$) to $11\%$
in San Francisco (where the chance rate is $4.2\%$). These detection
floors are tight enough that residual linear leakage of the magnitude
we should care about would be detected, and the test rejects no city
at the conventional $\alpha = 0.05$ threshold.

\paragraph{Downstream effect on the Baur DML estimate.} The Baur
PC1 estimates in Table~\ref{tab:baur-12city} can in principle be
recomputed on the post-LEACE+IGBP embedding to ask whether the
per-$\sigma$ effects we report would survive spatial deconfounding.
For the three markets in which we have both pre- and post-LEACE
estimates (Boston, New York, San Francisco, from the original
three-city paper version), the per-$\sigma$ PC1 effect drops in
magnitude after deconfounding but the New York effect remains
positive and significant. For the nine additional markets in
this revision we report the LEACE+IGBP guardedness verification but
not the post-deconfounding DML estimate, because the per-market
sample sizes after the LEACE residualization fall below the
threshold our pre-registered protocol requires for valid DML
inference. We flag this as a limitation in
Section~\ref{sec:conclusion} and note that subsequent work could
re-collect a larger per-market panel to address it directly.

\paragraph{Interpretation.} The LEACE+IGBP machinery developed here
establishes that the BERT representation in our analysis can be
cleanly deconfounded of lat-lon-targeted information at the linear
level (LEACE residual at numerical zero) and that nonlinear
residual leakage, where present, can be closed by IGBP within five
outer iterations. The Baur PC1 channel results in
Section~\ref{sec:baur-12city} are reported on the unprojected
BERT embedding, consistent with the published Baur, Rosenfelder, and
Lutz protocol, and the per-market sign and magnitude pattern we
document there should be read with the spatial-confounding caveat
that the deconfounding analysis here is set up to address. The
methodological contribution of this section is to demonstrate that
the linear concept erasure is exact in every metropolitan market we
consider, that the nonlinear residual closes within practical
compute, and that the binomial-power floors of the guardedness test
are tight enough to detect residual leakage at substantively
meaningful magnitudes.
